\chapter[Unified Geometric Interpretation]{Unified Geometric Interpretation}

The WIS framework reverses the traditional order of information
geometry.  Classical information geometry starts from a divergence and
derives a geometry (the Fisher metric from KL, for instance).  The WIS
framework first constructs a geometry---the pullback metric from the
logarithmic embedding---and then studies information-theoretic
quantities as scalar fields or quadratic forms on that fixed geometry.

This reversal has a structural consequence: every smooth information
functional, every Bregman divergence, and every piecewise-smooth
security measure obeys the same geometric stability estimates,
differing only in the constants that depend on their Hessian or
gradient norm.  The geometry is universal; the functionals are
individual.  The Universal Optimality Defect plays the role of the
canonical isotropic quadratic form with respect to which all other
divergences are anisotropic deformations, captured by the Average
Hessian Operator of Chapter~21.

Smooth functionals are controlled by gradient and Hessian estimates
(Theorems~14.1--14.2).  Bregman divergences are characterised through
the spectral decomposition of the Average Hessian Operator
(Theorem~15.3).  Rank-based security measures become smooth after
partitioning the posterior manifold into regularity intervals
(Chapter~23).  The analytical techniques differ, but the underlying
geometry is the same.